\begin{definition}
	Niech \(\mu, \nu\) będą rozkładami prawdopodobieństwa nad skończonym zbiorem \(S\). Sprzęganiem \(\mu\) i \(\nu\) nazywamy dowolną parę zmiennych losowych \(\left( X,Y \right) \) taką, że \(X\) ma rozkład \(\mu\), a \(Y\) ma rozkład \(\nu\). W szczególności te zmienne nie muszą być niezależne.
\end{definition}

\begin{lemma}
	Niech \(\left( X,Y \right) \) będzie sprzęganiem \(\mu\) i \(\nu\). Zachodzi
	\[ \left\|\mu-\nu\right\|_{TV} \le P\left( X \neq Y \right) . \]
	Ponadto istnieje sprzęganie dla którego zachodzi równość.
\end{lemma}
\begin{proof}
	Dla dowolnego \(A \subseteq S\) mamy
	\[ \mu\left( A  \right) - \nu\left( A  \right) = P\left( X \in A  \right) - P\left( Y \in  A  \right) \le P\left( X \in A \cap Y \notin A  \right) \le P \left( X \neq Y \right) . \]
	Analogicznie \(\nu\left( A  \right) - \mu\left( A  \right) \le  P\left( X \neq Y \right) \). To daje żądaną nierówność.

	Teraz skonstruujemy sprzęganie spełniające równość. Niech \(B = \left\{ x \in S : \mu\left( x  \right) \ge \nu\left( x  \right)  \right\} \). Niech \(p_1 = \mu\left( B  \right) - \nu\left( B  \right) \), \(p_2 = \nu\left( B^{c} \right) - \mu\left( B^{c} \right) \). Mamy \(p_1 = p_2 = \left\|\mu-\nu\right\|_{TV}\). Niech \(p_3 = 1-p_1 = 1-p_2\).

	Rzucamy monetą z prawdopodobieństwem orła \(p_3\). Jeśli wypadnie orzeł to ustalamy \(X = Y = s \), gdzie \(s\) wybieramy z \(S\) z rozkładem \(\left( \frac{1}{p_3} \min\left( \mu\left( s  \right) ,\nu\left( s  \right)  \right) : s \in S \right) \). Jeśli wypadnie reszka ustalamy \(X = x \) i \(Y=y\), gdzie \(x\) jest wybierany losowo z \(S \) z rozkładem \(\left( \frac{1}{p_1} \max\left( \mu\left( x  \right) -\nu\left( x  \right) , 0 \right) : x \in S \right) \), a \(y\) z rozkładem \(\left( \frac{1}{p_2} \max\left( \nu\left( x  \right) -\mu\left( x  \right) , 0 \right) : x \in S \right) \). W przypadku reszki jedna zmienna przyjmuje tylko te wartości, na których \(\mu\) jest większe, a druga tylko te, na których \(\nu\) jest większe. Mamy więc \(P\left( X\neq Y \right) = 1-p_3 = \left\|\mu-\nu\right\|_{TV}\), a \(\left( X,Y \right) \) faktycznie jest sprzęganiem \(\mu\) i \(\nu\) -- zmienne mają odpowiednie rozkłady.
\end{proof}